\documentclass[10pt,a4paper]{amsart}
\usepackage{amsmath,amssymb,amsthm,graphicx}
% Retained minimal preamble for the Stacks Project chapter varieties.tex.
% Verbatim excerpt of the chapter preamble: Xy-pic setup (native preamble lines 14-18), the native theorem
% declarations with the mathematical macros (native preamble lines 154-217), and the project own no-op
% definitions of the reference/slogan/history display environments (native wrapper.tex lines 1-3) that the
% counted statement uses for its slogan. Copied mechanically; nothing added, removed or redefined.
% For commutative diagrams we use Xy-pic
\usepackage[all]{xy}

% We use 2cell for 2-commutative diagrams.
\xyoption{2cell}
% Theorem environments.
%
\theoremstyle{plain}
\newtheorem{theorem}[subsection]{Theorem}
\newtheorem{proposition}[subsection]{Proposition}
\newtheorem{lemma}[subsection]{Lemma}

\theoremstyle{definition}
\newtheorem{definition}[subsection]{Definition}
\newtheorem{example}[subsection]{Example}
\newtheorem{exercise}[subsection]{Exercise}
\newtheorem{situation}[subsection]{Situation}

\theoremstyle{remark}
\newtheorem{remark}[subsection]{Remark}
\newtheorem{remarks}[subsection]{Remarks}

\numberwithin{equation}{subsection}

% Macros
%
\def\lim{\mathop{\mathrm{lim}}\nolimits}
\def\colim{\mathop{\mathrm{colim}}\nolimits}
\def\Spec{\mathop{\mathrm{Spec}}}
\def\Hom{\mathop{\mathrm{Hom}}\nolimits}
\def\Ext{\mathop{\mathrm{Ext}}\nolimits}
\def\SheafHom{\mathop{\mathcal{H}\!\mathit{om}}\nolimits}
\def\SheafExt{\mathop{\mathcal{E}\!\mathit{xt}}\nolimits}
\def\Sch{\mathit{Sch}}
\def\Mor{\mathop{\mathrm{Mor}}\nolimits}
\def\Ob{\mathop{\mathrm{Ob}}\nolimits}
\def\Sh{\mathop{\mathit{Sh}}\nolimits}
\def\NL{\mathop{N\!L}\nolimits}
\def\CH{\mathop{\mathrm{CH}}\nolimits}
\def\proetale{{pro\text{-}\acute{e}tale}}
\def\etale{{\acute{e}tale}}
\def\QCoh{\mathit{QCoh}}
\def\Ker{\mathop{\mathrm{Ker}}}
\def\Im{\mathop{\mathrm{Im}}}
\def\Coker{\mathop{\mathrm{Coker}}}
\def\Coim{\mathop{\mathrm{Coim}}}

% Boxtimes
%
\DeclareMathSymbol{\boxtimes}{\mathbin}{AMSa}{"02}

%
% Macros for moduli stacks/spaces
%
\def\QCohstack{\mathcal{QC}\!\mathit{oh}}
\def\Cohstack{\mathcal{C}\!\mathit{oh}}
\def\Spacesstack{\mathcal{S}\!\mathit{paces}}
\def\Quotfunctor{\mathrm{Quot}}
\def\Hilbfunctor{\mathrm{Hilb}}
\def\Curvesstack{\mathcal{C}\!\mathit{urves}}
\def\Polarizedstack{\mathcal{P}\!\mathit{olarized}}
\def\Complexesstack{\mathcal{C}\!\mathit{omplexes}}
% \Pic is the operator that assigns to X its picard group, usage \Pic(X)
% \Picardstack_{X/B} denotes the Picard stack of X over B
% \Picardfunctor_{X/B} denotes the Picard functor of X over B
\def\Pic{\mathop{\mathrm{Pic}}\nolimits}
\def\Picardstack{\mathcal{P}\!\mathit{ic}}
\def\Picardfunctor{\mathrm{Pic}}
\def\Deformationcategory{\mathcal{D}\!\mathit{ef}}
\newenvironment{reference}{}{}
\newenvironment{slogan}{}{}
\newenvironment{history}{}{}
\usepackage{cleveref}
\crefname{originalreference}{original reference}{original references}
\Crefname{originalreference}{Original reference}{Original references}
\makeatletter
\newlabel{constructions-equation-multiply-on-sheaf}{{03GJ}{source}{}{sourcepointer-1}{}}
\newlabel{constructions-equation-multiply-on-sheaf@cref}{{[originalreference][0][]03GJ}{[1][0][]source}{}{}{}}
\makeatother
\crefname{originalreference}{original reference}{original references}
\Crefname{originalreference}{Original reference}{Original references}
\makeatletter
\newlabel{algebra-lemma-ass-zero-divisors}{{00LD}{source}{}{sourcepointer-2}{}}
\newlabel{algebra-lemma-ass-zero-divisors@cref}{{[originalreference][0][]00LD}{[1][0][]source}{}{}{}}
\makeatother
\crefname{originalreference}{original reference}{original references}
\Crefname{originalreference}{Original reference}{Original references}
\makeatletter
\newlabel{algebra-lemma-finite-ass}{{00LC}{source}{}{sourcepointer-3}{}}
\newlabel{algebra-lemma-finite-ass@cref}{{[originalreference][0][]00LC}{[1][0][]source}{}{}{}}
\makeatother
\crefname{originalreference}{original reference}{original references}
\Crefname{originalreference}{Original reference}{Original references}
\makeatletter
\newlabel{divisors-definition-zero-scheme-s}{{02OQ}{source}{}{sourcepointer-4}{}}
\newlabel{divisors-definition-zero-scheme-s@cref}{{[originalreference][0][]02OQ}{[1][0][]source}{}{}{}}
\makeatother
\crefname{originalreference}{original reference}{original references}
\Crefname{originalreference}{Original reference}{Original references}
\makeatletter
\newlabel{coherent-lemma-i-star-equivalence}{{087T}{source}{}{sourcepointer-5}{}}
\newlabel{coherent-lemma-i-star-equivalence@cref}{{[originalreference][0][]087T}{[1][0][]source}{}{}{}}
\makeatother
\title{Existence of a section of O(1) giving an injective multiplication F(-1) -> F with cokernel pushed forward from the associated hyperplane}
\author{The Stacks Project Authors}
\date{stacks-project:varieties.tex:local-snapshot-20261011; GNU FDL 1.2 or later}
\begin{document}
\maketitle
\noindent\textit{Editorial note.} All delivered mathematical text is the author's own, copied verbatim from the pinned Stacks Project chapter file varieties.tex; no statement or proof step was rewritten, reordered or removed.\par\medskip
\noindent\textit{Editorial note.} The counted span cites five cross-chapter labels (including an equation and a definition), each recorded in assembly.external\_references with its canonical tag as a navigation pointer.\par\medskip
\noindent\textit{Editorial note.} The proof cites no bibliography item inside the counted span.\par\medskip
\noindent\textit{Allowed standard prerequisites (editorial).}\begin{itemize}
\item Associated primes of a finite module over a Noetherian ring and the zero-divisor criterion. Stacks Project, Algebra, Lemmas algebra-lemma-finite-ass (tag 00LC) and algebra-lemma-ass-zero-divisors (tag 00LD). Finiteness of Ass(M) and the fact that an element outside all associated primes is a non-zero-divisor drive the injectivity argument.
\item Multiplication by a section on a quasi-coherent sheaf. Stacks Project, Constructions, Equation constructions-equation-multiply-on-sheaf (tag 03GJ). Defines exactly the map F(-1) -> F used in the statement.
\item The zero scheme of a section of an invertible sheaf and i\_*-equivalence on the hyperplane. Stacks Project, Divisors, Definition divisors-definition-zero-scheme-s (tag 02OQ) and Cohomology of Schemes, Lemma coherent-lemma-i-star-equivalence (tag 087T). Identify the ideal sheaf of H and descend the cokernel to the hyperplane.
\end{itemize}\medskip
\ifdefined\hypertarget\hypertarget{sourcepointer-1}{}\fi
\noindent\textit{Source reference.} Original reference 03GJ: sources/stacks-chow/tags.txt, source lines 4530--4530. Cross-chapter standard prerequisite cited by the author inside the counted span; it is used as a pointer and its own proof is not part of this item.\par\medskip
\ifdefined\hypertarget\hypertarget{sourcepointer-2}{}\fi
\noindent\textit{Source reference.} Original reference 00LD: sources/stacks-chow/tags.txt, source lines 783--783. Cross-chapter standard prerequisite cited by the author inside the counted span; it is used as a pointer and its own proof is not part of this item.\par\medskip
\ifdefined\hypertarget\hypertarget{sourcepointer-3}{}\fi
\noindent\textit{Source reference.} Original reference 00LC: sources/stacks-chow/tags.txt, source lines 782--782. Cross-chapter standard prerequisite cited by the author inside the counted span; it is used as a pointer and its own proof is not part of this item.\par\medskip
\ifdefined\hypertarget\hypertarget{sourcepointer-4}{}\fi
\noindent\textit{Source reference.} Original reference 02OQ: sources/stacks-chow/tags.txt, source lines 3525--3525. Cross-chapter standard prerequisite cited by the author inside the counted span; it is used as a pointer and its own proof is not part of this item.\par\medskip
\ifdefined\hypertarget\hypertarget{sourcepointer-5}{}\fi
\noindent\textit{Source reference.} Original reference 087T: sources/stacks-chow/tags.txt, source lines 10438--10438. Cross-chapter standard prerequisite cited by the author inside the counted span; it is used as a pointer and its own proof is not part of this item.\par\medskip
\begin{lemma}
\label{lemma-exact-sequence-induction}
Let $k$ be an infinite field. Let $n \geq 1$.
Let $\mathcal{F}$ be a coherent module on $\mathbf{P}^n_k$.
Then there exist a nonzero section
$s \in \Gamma(\mathbf{P}^n_k, \mathcal{O}(1))$
and a short exact sequence
$$
0 \to \mathcal{F}(-1) \to \mathcal{F} \to i_*\mathcal{G} \to 0
$$
where $i : H \to \mathbf{P}^n_k$ is the hyperplane $H$ associated to $s$
and $\mathcal{G} = i^*\mathcal{F}$.
\end{lemma}

\begin{proof}
The map $\mathcal{F}(-1) \to \mathcal{F}$ comes from
Constructions, Equation (\ref{constructions-equation-multiply-on-sheaf})
with $n = 1$, $m = -1$ and the section $s$ of $\mathcal{O}(1)$.
Let's work out what this map looks like if we restrict it to
$D_{+}(T_0)$. Write $D_{+}(T_0) = \Spec(k[x_1, \ldots, x_n])$
with $x_i = T_i/T_0$. Identify $\mathcal{O}(1)|_{D_{+}(T_0)}$ with
$\mathcal{O}$ using the section $T_0$. Hence if
$s = \sum a_iT_i$ then $s|_{D_{+}(T_0)} = a_0 + \sum a_ix_i$
with the identification chosen above. Furthermore, suppose
$\mathcal{F}|_{D_{+}(T_0)}$ corresponds to the finite
$k[x_1, \ldots, x_n]$-module $M$. Via the identification
$\mathcal{F}(-1) = \mathcal{F} \otimes \mathcal{O}(-1)$
and our chosen trivialization of $\mathcal{O}(1)$ we see that
$\mathcal{F}(-1)$ corresponds to $M$ as well. Thus
restricting $\mathcal{F}(-1) \to \mathcal{F}$ to $D_{+}(T_0)$ gives
the map
$$
M \xrightarrow{a_0 + \sum a_ix_i} M
$$
To see that the arrow is injective, it suffices to pick
$a_0 + \sum a_ix_i$ outside any of the associated primes of $M$, see
Algebra, Lemma \ref{algebra-lemma-ass-zero-divisors}. By
Algebra, Lemma \ref{algebra-lemma-finite-ass}
the set $\text{Ass}(M)$ of associated primes of $M$ is finite.
Note that for $\mathfrak p \in \text{Ass}(M)$ the intersection
$\mathfrak p \cap \{a_0 + \sum a_i x_i\}$ is a proper $k$-subvector space.
We conclude that there is a finite family of proper sub vector spaces
$V_1, \ldots, V_m \subset \Gamma(\mathbf{P}^n_k, \mathcal{O}(1))$
such that if we take $s$ outside of $\bigcup V_i$, then multiplication
by $s$ is injective over $D_{+}(T_0)$.
Similarly for the restriction to $D_{+}(T_j)$ for $j = 1, \ldots, n$.
Since $k$ is infinite, a finite union of proper sub vector spaces
is never equal to the whole space, hence we may choose $s$ such
that the map is injective.
The cokernel of $\mathcal{F}(-1) \to \mathcal{F}$ is annihilated
by $\Im(s : \mathcal{O}(-1) \to \mathcal{O})$ which is
the ideal sheaf of $H$ by
Divisors, Definition \ref{divisors-definition-zero-scheme-s}.
Hence we obtain $\mathcal{G}$ on $H$ using
Cohomology of Schemes, Lemma \ref{coherent-lemma-i-star-equivalence}.
\end{proof}

\end{document}
